\subsection{The fixed-field rank-four trace integer}
\label{subsec:sr-rank-four-trace-integer}
We make explicit that the large auxiliary CM field used to
produce support is absent from the numerical estimate.
Only the actual descended representations and their original
local parameters enter this subsection.

Fix eighteen imaginary quadratic fields \(K_i\) with
pairwise coprime prime discriminants.
They are fixed once and for all.
For any rational prime \(p\), at most one discriminant
is divisible by \(p\).
Lemma~\ref{lem:sr-quadratic-self-twists} shows that
a rank-four irreducible representation has reducible
restriction to at most sixteen such quadratic fields.
Thus for every irreducible rational rank-four point
there is a \(K_i\) such that
\[
           p\nmid\operatorname{disc}K_i,\qquad
                         r|_{G_{K_i}}\text{ is irreducible}.
 \tag{Q1}
\]
For an irreducible fixed restriction there are at most
two extensions to \(G_{\mathbb Q}\), by Schur's lemma
and the quadratic character of \(K_i/\mathbb Q\).

\begin{srproposition}\label{prop:sr-rank-four-fixed-integer}
Use the actual Case II support construction and
Proposition~\ref{prop:sr-four-actual-final-endpoint},
including the automorphy and local conclusions proved in
Section~\ref{subsec:sr-case-two-descent}.
For its admissible rational points, the rank-four
common-integer comparison of
Proposition~\ref{prop:sr-rank-four-absorption}
holds with constants \(C_4,A_4\) fixed before
\(M,a,p,k\), and before the coefficient prime and field.
\end{srproposition}
\begin{proof}
The arithmetic weights of the original rational point are
\((a+1,a,1,0)\), \(a\ge2\).
They are paired with common sum \(a+1\).
After restriction to \(K_i\), their regular pure infinity
parameter is cohomological for an algebraic coefficient
system whose weight sizes are bounded by a fixed multiple
of \(a+1\).
This is the same cohomological criterion and highest-weight
construction used in Section~\ref{sec:spectrum-deformation}.
The criterion applies to every regular algebraic cusp
over a totally imaginary field: its pure infinity
parameter has nonzero relative Lie algebra cohomology.
It does not require the finite polarization character
to be trivial.\footnote{L.~Clozel,
\href{https://www.jmilne.org/math/Books/AA1988a.pdf}
{\emph{Motifs et formes automorphes: applications du
principe de fonctorialit\'e}}, Lemma 3.14 (complex places),
with the purity lemma used there; Theorem 3.13 gives
the coefficient conjugation assertion used below.}
Equivalently one may count all cohomological cusps with
these infinity parameters and retain their polarized subset.

We verify the level on the descended field.
The actual point has ordinary rank-one local graded
characters, finite inertia away from the coefficient prime,
and zero Weil--Deligne monodromy.
At that prime the local no-edge computation and
Proposition~\ref{prop:sr-actual-wit} give the same
character parameter with zero monodromy.
Every finite inertial character is a restriction of a
character whose conductor divides the original level \(M\).
On restriction to \(K_{i,v}\), its conductor divides
\(M\mathcal O_{K_{i,v}}\):
local norms of units congruent to one modulo this ideal
are congruent to one modulo \(M\) over the rational
completion.
Full local--global compatibility identifies the actual
automorphic local component with the corresponding
unlinked principal series.
The induced-function vector is fixed by the principal
congruence subgroup of level \(M\).
Consequently the descended cusp has a \(K_4(M)\)-fixed
vector over \(K_i\).
This argument is on the original point and its finite
characters; no level of the auxiliary field \(L\) occurs.

The proved coefficient/level cohomology bound therefore
gives, for the collection \(\Sigma_i(M,a)\) of all such
global cuspidal Hecke systems, fixed constants
\[
 |\Sigma_i(M,a)|\le C_i[M(a+1)]^{A_i}.
\]
Take the sum of the constants and the maximum of these
eighteen fixed exponents.
After enlarging one fixed constant we obtain
\[
                \sum_i|\Sigma_i(M,a)|
                       \le C_0[M(a+1)]^{A_0}.
 \tag{Q2}
\]
Each collection is algebraic and stable under coefficient
conjugation by Clozel's Theorem 3.13.
That theorem supplies cuspidal conjugates, rather than
assuming that the cuspidal summand of Betti cohomology
already has a rational structure.
The compact congruence subgroup is defined over
\(\mathbb Q\), so coefficient conjugation preserves
its fixed vectors, and the parallel infinity parameters
are unchanged.
Its cardinality and all constants in (Q2) are independent
of the field used in the support construction.

Fix \(p\ge5\), \(p\nmid M\).
For each \(i\) with \(p\nmid\operatorname{disc}K_i\), put
\(f_i=1\) or \(2\), its residue degree at \(p\).
A system in \(\Sigma_i\) supplies the four arithmetic
Frobenius eigenvalues \(\beta_1,\ldots,\beta_4\) over
the residue field of size \(p^{f_i}\).
Form every multiset
\[
                 \{\alpha_1,\ldots,\alpha_4\},
                     \qquad \alpha_j^{f_i}=\beta_j,
\]
and its trace \(t=\sum_j\alpha_j\).
There are at most \(2^4=16\) choices.
Retain two labelled copies of every choice.
This defines a finite Galois-stable trace list \(S\) with
\[
                       |S|\le32C_0[M(a+1)]^{A_0}.
 \tag{Q3}
\]
Labels with equal numerical trace are retained.

This list classifies every finite pairwise nonisomorphic
family of actual rational points in the sense needed by
Section~\ref{sec:common-trace-integer}.
Assign each point one field satisfying (Q1), and then
its descended global Hecke system there.
The rational Frobenius eigenvalues are among the indicated
\(f_i\)-th roots.
For a given descended system at most two rational points
occur, by the extension-fiber assertion above.
The two labelled copies therefore make the label
assignment injective on that finite family.
No assertion of distinct traces is needed.

For completeness all these candidate traces are algebraic
integers in the present arithmetic normalization.
At coefficient primes away from \(p\), continuous
compatible Galois representations have stable lattices,
giving integrality there.
At \(p\), the local automorphic component is unramified,
so its Galois representation is crystalline.
The original arithmetic Hodge numbers have smallest
exponent zero.
In our convention the filtration indices of the
crystalline module are the negatives of these exponents,
and arithmetic Frobenius corresponds to inverse crystalline
Frobenius.
The upper Newton-polygon bound from weak admissibility
therefore makes those inverse eigenvalues integral
at \(p\).
This is the integral Satake estimate used in
Section~\ref{sec:common-trace-integer}, with its
normalization made explicit.
Thus the \(\beta_j\) are integral at every finite place,
and every \(\alpha_j\), being a root of
\(X^{f_i}-\beta_j\), is an algebraic integer.
The compatible realizations here are those of the
already polarized regular algebraic cusps.
No automorphy lifting theorem at a small coefficient
prime is being invoked.

We next prove separation.
The single global unitarizing norm twist has center
\((a+1)/2\), determined by the central character.
At arithmetic Frobenius the generic unitary bound gives
\[
       p^{a/2}<|\alpha_j|<p^{(a+2)/2},\qquad
                       \prod_j|\alpha_j|=p^{2(a+1)} .
\]
The same inequalities hold for every candidate root,
since its absolute value is determined by
\(\alpha_j^{f_i}=\beta_j\).
Put \(x_j=|\alpha_j|/p^{(a+1)/2}\).
Then \(p^{-1/2}\le x_j\le p^{1/2}\) and \(\prod x_j=1\).
On the compact polytope
\(\sum\log x_j=0\), \(|\log x_j|\le\frac12\log p\),
the convex function \(\sum e^{\log x_j}\) attains its
maximum at a vertex.
A vertex has two coordinates at each endpoint.
Consequently
\[
                  |t|\le
                       2\bigl(p^{(a+2)/2}+p^{a/2}\bigr).
 \tag{Q4}
\]
For any roots of unity \(z_1,z_2,z_3\),
\[
 |p^{a+1}z_1+p^az_2+pz_3+1|
                       \ge p^{a+1}-p^a-p-1
                       >2\bigl(p^{(a+2)/2}+p^{a/2}\bigr).
 \tag{Q5}
\]
To check the last inequality, divide by \(p^{a/2}\).
Its left-minus-right increases with \(a\ge2\).
At \(a=2\) it is
\(p^2-3p-3-p^{-1}>0\) for \(p\ge5\).
This proves every trace difference nonzero.

Let \(\mathcal Z\) be the Galois orbit of the prescribed
triple \((z_1,z_2,z_3)\), with all entries of order
dividing \(k\).
It has at most \(\varphi(k)\) elements.
Define, before choosing a coefficient realization,
\[
 J=\prod_{z\in\mathcal Z}\prod_{s\in S}
       \bigl(t_s-(p^{a+1}z_1+p^az_2+pz_3+1)\bigr).
 \tag{Q6}
\]
Every factor is integral and nonzero.
Galois stability makes \(J\) a nonzero rational
algebraic integer, hence a nonzero integer.
As in the common trace product proof, the injective label
assignment and integrality permit discarding unused factors
at every \(\ell\ne p\).
Thus
\[
                       |J|_\ell\le\prod_r|d_r|_\ell.
\]
The uniform archimedean bound for a factor is
\(8p^{a+1}\), whose logarithm is at most
\(2(a+1)\log p\).
Equations (Q3) and (Q6) give
\[
 \log|J|\le64C_0\varphi(k)
                    [M(a+1)]^{A_0}(a+1)\log p
           \le C_4\varphi(k)[M(a+1)]^{A_4}\log p,
\]
where \(C_4=64C_0\) and \(A_4=A_0+1\).
This accounts for the weight factor inside the final
fixed exponent; no extra \(a+1\) is left outside it.
The quantifier order and normalization are exactly those
of \(\mathrm{NumberSp4Alt}\).
\end{proof}

\subsection{Distinct marked Gaussian points after restriction}
\label{subsec:sr-gaussian-restriction-classes}
\begin{srlemma}\label{lem:sr-gaussian-marked-restriction-injective}
Let two actual integral normalized rank-three unitary
type-I points have isomorphic restrictions to \(G_K\).
Then they differ by the marked torus scaling action.
In particular distinct mixed torus orbits give distinct
irreducible \(G_K\)-representations.
\end{srlemma}
\begin{proof}
Write \(\theta(A)=J_3A^{-\mathsf t}J_3^{-1}\).
Let \(A\rho_1A^{-1}=\rho_2\).
The two unitary laws show that
\(\theta(A)^{-1}A\) commutes with \(\rho_1(G_K)\).
It is scalar by irreducibility.
Thus \(A\theta(A)^{-1}=b1\) for some nonzero \(b\).
Replacing \(A\) by \(\lambda A\), with
\(\lambda^2b=1\), leaves its conjugation on \(G_K\)
unchanged and makes its semidirect conjugation fix
the normalized value at \(c\).
It now intertwines the full normalized semidirect points.

For either point the marked difference at \(g_1\) is
\(\operatorname{diag}(\zeta_i,0,-\zeta_i)\).
The intertwiner makes these differences conjugate.
Their nonzero eigenvalues give
\(\zeta_2=\zeta_1\) or \(-\zeta_1\).
Both \(\zeta_i\) are integral and reduce to the same
nonzero prescribed value.
The negative alternative would give \(2\bar\zeta_1=0\),
impossible at odd \(\ell\).
Hence \(\zeta_2=\zeta_1\), and the intertwiner is diagonal.
The relation \(A=\theta(A)\) gives
\(A=\operatorname{diag}(u,b,u^{-1})\) with \(b^2=1\).
Multiplying by \(-1\) if needed, which does not change
ordinary matrix conjugation, makes \(b=1\).
The resulting intertwiner is
\(\operatorname{diag}(t,1,t^{-1})\), the actual marked
scaling action.
\end{proof}

\subsection{The uniform nontrivial exceptional and resonant estimate}
\label{subsec:sr-uniform-nontrivial-assembly}
The Case II rank-four support and endpoint are
Proposition~\ref{prop:sr-four-actual-final-endpoint}.
The Gaussian two-run output is
Proposition~\ref{prop:sr-gaussian-actual-g10}.
We assemble their estimates, recording where the
auxiliary degree disappears.

\begin{srproposition}\label{prop:sr-nontrivial-final-assembly}
There is a single natural exponent \(B\), fixed
before the principal host, \(a,\ell,\psi\) and the splitting
ring, such that
\[
                  (\log\ell)v_{\ell,a}(\psi)
                                      \le(\ell Na)^B
\]
for every nontrivial principal character in the required
exceptional or resonant packet.
\end{srproposition}
\begin{proof}
The actual host is the fixed field in
Section~\ref{sec:sr-arithmetic}, with its even dyadic
character packet.
For \(a\ge2\),
Proposition~\ref{prop:sr-principal-case-II-arithmetic}
puts every required instance in Case II, with
\(\ell\nmid N\).
The actual seeds, local witnesses, Hecke kernels and
selected-tower irreducibility have been matched to the
shared common-module construction.
Proposition~\ref{prop:sr-actual-corx} supplies the
residual unramified-cocycle vanishing needed by the
arithmetic data constructor and the marked comparison.
Proposition~\ref{prop:sr-four-actual-final-endpoint}
and the exact descent give \(\mathrm{IsAutQA}\) and the
original local parameters.
The residual no-quadratic-self-twist argument gives
cyclotomic-kernel irreducibility.
The ordinary tangent calculation places the tangent
classes in the descended polarized adjoint
Bloch--Kato group, whose vanishing gives rigidity.
Here the base of that adjoint group is the totally
real field, as required by the polarized vanishing
theorem; vanishing of the full adjoint group over
a CM field is not asserted.

The marked cotangent injection and the two Euler/Rees
comparisons then give the normalized trace inequality
\[
 (\log\ell)v_{\ell,a}(\psi)
                     \le2\sum_r-\log|d_r|_\ell.
 \tag{Q7}
\]
These are the same actual finite representatives and
the same local deformation problem as in the count.
Proposition~\ref{prop:sr-rank-four-fixed-integer}
supplies their common integer on fixed fields at
the original level \(M=\ell2^nN\).
Proposition~\ref{prop:sr-rank-four-absorption} gives
\[
 B_4=10A_4+12+b_4,\qquad 2C_4\le2^{b_4}.
\]
All these constants are fixed.
In particular no estimate polynomial in \([L:\mathbb Q]\)
or in the Ferrero--Washington constant is used:
\(L\), its large local degrees and the seed weight
were used only to obtain support, and descent preceded
(Q2).

For \(a=1\), the required primes are \(3,7\).
The small dyadic-conductor branches satisfy
Proposition~\ref{prop:sr-small-conductor-exceptional},
with
\[
 B_{\rm small}=196609\cdot393216^{24}.
\]
In the large branches the actual Gaussian construction
and Proposition~\ref{prop:sr-gaussian-actual-g10}
give finiteness,
strong-open irreducibility and exact automorphy over
the fixed field \(\mathbb Q(i)\).
Proposition~\ref{prop:sr-gaussian-rigidity}
and Lemma~\ref{lem:sr-gaussian-actual-level} give
rigidity, crystalline local behavior and conductor
norm at most \(N^4\), at the fixed weight
\((0,-1,-2)\).
Proposition~\ref{prop:sr-actual-corx} supplies its
residual corner-vanishing condition.
Lemma~\ref{lem:sr-gaussian-marked-restriction-injective}
checks the distinctness needed by the trace product.
Counting is therefore on this fixed Gaussian spectrum,
rather than on the growing auxiliary field.
At \(p=11\),
Proposition~\ref{prop:sr-gaussian-fixed-prime} yields
\[
 B_G=A_G+1+b_G,\qquad
                      4C_G2^{A_G}\log11\le3^{b_G}.
\]
The two possible marked embeddings are exactly the
factor two already retained in that proposition.

Set \(B=\max(B_4,B_G,B_{\rm small})\).
This choice precedes every principal datum.
The normalized strict length is invariant under finite
splitting-coefficient extension, as already proved,
so the estimate holds in every permitted coefficient ring.
The three exhaustive branches prove the assertion.
\end{proof}

\paragraph{Exact producer scope.}
The rank-four ideal factorization, finite weight zero,
no-full-block and all-irreducible-point statements are
proved in Proposition~\ref{prop:sr-four-actual-final-endpoint}.
The actual Gaussian frame, common coefficient modules,
two seeds, rank-two run, outer-corner ideal, rank-three
run and all-dominant classical realization are instantiated
in Proposition~\ref{prop:sr-gaussian-actual-g10}.
Thus the support outputs formerly denoted (G10) are
conclusions of those constructions.
The subsequent finite-part, exact descent, original-level,
cyclotomic-irreducibility, local tangent, fixed-spectrum,
integer-product and scalar-absorption steps leave no
auxiliary-degree dependence in \(B\).
